\section{Step One: Verification of DNS Role
Installation}\label{step-one-verification-of-dns-role-installation}

1.1: On your server, open Server Manager.

1.2: Click on Tools in the top-right corner and select DNS to open the
DNS Manager.

S1.3: Verify that your server (GROUP1-SVR) appears in the console tree.

\includegraphics[width=6.01088in,height=4.12432in]{media/image1.png}\textbf{\hfill\break
\hfill\break
}

\section{Step 2: Configuring the Forward Lookup
Zone}\label{step-2-configuring-the-forward-lookup-zone}

The Forward Lookup Zone is responsible for resolving hostnames (like
your server name) into IP addresses.

2.1. In the DNS Manager, expand the node for GROUP1-SVR.

2.2. Expand Forward Lookup Zones.

2.3. Verify that the primary zone for your domain i.e. Group1.local was
created during the AD DS setup.

2.4: Click on the zone to view the automatically generated records (SOA,
NS, and A records).

\includegraphics[width=6in,height=4.82292in]{media/image2.png}

\textbf{\hfill\break
}

\textbf{\hfill\break
}

\textbf{Step 3: Configuring the Reverse Lookup Zone}

The Reverse Lookup Zone resolves IP addresses back into hostnames, which
plays a pivotal role in network troubleshooting and certain app
authentications.

3.1. Right-click on Reverse Lookup Zones in the DNS Manager and select
New Zone.

3.2. Click Next on the Welcome screen.

3.3. Select Primary zone and ensure "Store the zone in Active Directory"
is checked. Click Next.

3.4. Choose to replicate data to all DNS servers running on domain
controllers in this domain.

3.5. Select IPv4 Reverse Lookup Zone and click Next.

3.6. Enter your Network ID (i.e. 192.168.64.2) that corresponds to your
Windows Server\textquotesingle s static IP. Click Next.

3.7. Allow only secure dynamic updates (recommended for Active
Directory) and click Finish.

\includegraphics[width=6in,height=4.20833in]{media/image3.png}

\section{Step 4: Adding and Verifying DNS
Records}\label{step-4-adding-and-verifying-dns-records}

4.1. Navigate back to your Forward Lookup Zone (Group1cput.local).

4.2. Locate the Host (A) record for Group {[}X{]} Windows Server 2019.
Double-click it to open its properties.

4.3. Click the box that says, "Update associated pointer (PTR) record"
and click Apply and OK.

4.4. Navigate to your newly created Reverse Lookup Zone and refresh.
Verify that a PTR record pointing to your server now exists.

\textbf{\hfill\break
}\includegraphics[width=6in,height=4.72292in]{media/image4.png}

\section{5. Testing DNS Resolution (Server and
Client)}\label{testing-dns-resolution-server-and-client}

Testing confirms that the DNS server can successfully resolve names and
IP addresses, which is necessary to achieve full marks for the testing
requirement.

5.1 (Server-side test): Open PowerShell (or Command Prompt if
restrictions aren\textquotesingle t applied to admins).

Step 5.2: Type nslookup and press Enter. The default server should
reflect your server\textquotesingle s name and IP address.

Step 5.3: Perform a forward lookup by typing your
server\textquotesingle s fully qualified domain name (FQDN) (e.g.
Group1WindowsServer2019.Group1cput.local) and pressing Enter. It should
return your server\textquotesingle s IP address.

Step 5.4: Perform a reverse lookup by typing your
server\textquotesingle s IP address (e.g., 192.168.64.2) and press
Enter. It should return the FQDN.

\includegraphics[width=6in,height=4.68194in]{media/image5.png}

Step 5.5 (Client-side test): Switch to a client virtual machine running
on your Semester 1 hypervisor that is connected to the same virtual
network.

Step 5.6: Open Command Prompt/PowerShell on the client and type ping
Group1cput.local or ping the server by name (ping
Group1WindowsServer2019).

Step 5.7: Verify that the client successfully resolves the name to the
server\textquotesingle s IP address and receives replies.

\textbf{\hfill\break
\hfill\break
}\includegraphics[width=6in,height=4.76944in]{media/image6.png}

\textbf{\hfill\break
}

\subsection{References}\label{references}

McDonald,~G.,~Papadopoulos,~P.,~Pitropakis,~N.,~Ahmad,~J.,~\&
Buchanan,~W.~J.~(2022).~Ransomware:~Analysing the impact on Windows
Active Directory Domain
Services.~\emph{Sensors},~\emph{22},~953.~\url{https://doi.org/10.3390/s22030953}~Cited
by:~69

Mokhtar,~B.,~Jurcut,~A.,~ElSayed,~M.,~\& Azer,~M.~(2022).~Active
Directory attacks---Steps,~types,~and
signatures.\emph{Electronics},~\emph{11},~2629.~\url{https://doi.org/10.3390/electronics11162629}~Cited
by:~41

Moric,~Z.,~Dakic,~V.,~Redzepagic,~J.,~\& Otocan,~F.~(2024).~Penetration
testing and exploitation of Active Directory configuration
vulnerabilities.~\emph{2024 4th Intelligent Cybersecurity Conference
(ICSC)},~118--126.\url{https://doi.org/10.1109/icsc63108.2024.10895772}~Cited
by:~4

Pongsrisomchai,~S.,~\& Ngamsuriyaroj,~S.~(2019).~Automated IT audit of
Windows Server access control.~\emph{2019 21st International Conference
on Advanced Communication Technology
(ICACT)},~539--544.\url{https://doi.org/10.23919/icact.2019.8701931}~Cited
by:~1
